\section{Related work}
\paragraph{Memory and applicability.}
Voyager stores reusable executable skills, while ExpeL extracts natural language lessons from experience without updating model parameters~\citep{voyager,expel}. These systems motivate procedural reuse, but successful transfer is conditional on what the memory preserves. \citet{mtl} study four memory representations across six coding benchmarks and find that abstract insights transfer more reliably than concrete traces. Our experiment fixes a procedure valid at the source and changes its relevance to a target security contract; the endpoint distinguishes failure to complete from functional completion that fails a focal witness.

Boundary-Aware Skill Memory (BASM) is the closest applicability aware method~\citep{basm}. It augments skills with applicability conditions, risk cues, avoidance rules, and recovery notes, and combines retrieval, conditioned execution, and local repair. It also investigates wrong tool preference using attention and intervention analyses. Our B arm is a single generic reminder attached to a supplied memory. It neither implements BASM nor tests its structured mechanism. Prior work asks whether skills transfer or how to make them aware of applicability boundaries. We hold a source valid procedure fixed, change the target assumption, and separately observe functionality and a focal security property. We do not claim novelty for the general observation that a relevant skill can become inapplicable.

\paragraph{Skill induced regressions.}
SWE-Skills-Bench measures the marginal utility of supplied skills in repository software tasks and reports limited average gains and failures of context compatibility~\citep{sweskills}. \citet{harmfulskills} use differential analysis of paired runs to examine functional failures and efficiency regressions induced by loaded skills. Their results make clear that procedural guidance can constrain implementation in unhelpful ways. Our focal distinction is between failing to implement a feature and implementing it while violating a separately evaluated focal property. Our behavioral analysis separates visible statements, implementation changes, and measured outcomes.

\paragraph{Security and the execution interface.}
SecureVibeBench reconstructs vulnerability introducing scenarios in 105 C/C++ tasks from 41 OSS-Fuzz projects and combines functionality and security evaluation~\citep{securevibebench}. We cite its ACL 2026 version, which supersedes the earlier arXiv report. Its repository realism addresses a different objective from our synthetic memory intervention. Experience driven self evolution can also trade safety against task utility in web and embodied settings~\citep{experience_safety}; our measured behavior is the generated implementation rather than refusal of a harmful request. Finally, SWE-agent demonstrates that the agent computer interface affects software engineering performance~\citep{sweagent}. This matters directly here: different tool surfaces and execution budgets are part of each evaluated configuration, and interface related failure to complete changes the meaning of a focal witness pass.
